\section{Verification Contracts and Defenses}
\label{sec:defenses}
\label{sec:verification-contracts}

\subsection{State-and-secrecy contracts}

ReQWIRE and later uncomputation work show how useful it is to prove that
ancillas return to an asserted state \cite{rand2019reqwire,venev2024uncomputation}.
We propose extending the contract from a state assertion to a security
assertion.

A subroutine contract should contain:

\begin{enumerate}
  \item \textbf{Precondition}: initialization and allowed input subspace.
  \item \textbf{Postcondition}: exact returned state of each released qubit.
  \item \textbf{Separability}: released qubits factor from all live registers.
  \item \textbf{Noninterference}: the reduced state of every attacker-visible
  register is independent of protected variables.
  \item \textbf{Liveness}: no authorized future operation uses the old logical
  identity.
\end{enumerate}

An illustrative interface is:

\begin{lstlisting}[style=python]
@quantum_contract(
    requires={"ancilla": "|0>", "destination": "|0>"},
    ensures={
        "source": "|0>",
        "ancilla": "|0>",
        "destination": "input_state(source)",
        "separable": ["source", "ancilla"],
    },
    forbids_flow=[("secret", "released_or_hidden_qubits")],
)
def coherent_state_transfer(source, ancilla, destination):
    ...
\end{lstlisting}

\subsection{Complementary-basis testing}

A system that validates only computational-basis outputs is vulnerable to phase
and coherence changes. Critical kernels should include randomized tests in
complementary bases, entangled-reference tests, or process-level equivalence
checking. For Clifford subcircuits, stabilizer methods can often verify these
properties efficiently. More general circuits may require symbolic methods,
decision diagrams, randomized identities, or bounded formal verification.

\subsection{Cross-layer attestation}

Security should bind together:

\begin{itemize}
  \item the user-level circuit;
  \item the transpiled native-gate circuit;
  \item logical-to-physical qubit mapping;
  \item timing and scheduling;
  \item pulse or analog control definitions;
  \item measurement destinations and returned classical data.
\end{itemize}

Hashing only the high-level circuit is inadequate if lower layers can reorder
qubits or alter pulse semantics. Pulse-level attacks demonstrate why semantic
verification must extend below gate names \cite{xu2024pulseattacks}.

\subsection{Trap and canary circuits}

A provider or client can interleave secret trap circuits whose correct behavior
is difficult for an adversary to distinguish from real workloads. Useful traps
would include phase-sensitive tests, randomized Clifford identities, known
entangled states, and circuits that expose unauthorized coupling to nominally
idle qubits.

\subsection{Isolation and least privilege}

Source controls, basis choices, key-dependent classical branches, and calibration
interfaces should be treated as secrets. Compiler passes should receive only the
information they need. Multi-tenant devices should isolate qubit neighborhoods,
control channels, and readout paths to reduce crosstalk and unauthorized
observation.
